% PreTextuais/Abstract.tex
\begin{resumo}[ABSTRACT]
\begin{otherlanguage*}{english}
Public fleet management has historically operated under reactive emergency routines due to fragmented administrative records across asset inventories, maintenance work orders (WOs), and fuel logs. Guided by the logical method for scientific writing (Volpato, 2015, 2017), this master's thesis formulates and benchmarks a prescriptive analytics architecture driven by Explainable Artificial Intelligence (XAI), Graph Theory, and Mixed-Integer Linear Programming (MILP), prioritizing conceptual rigor over mechanical templates and categorially separating statistical significance from practical operational relevance. The empirical foundation comprises 77,867 work orders executed between 2020 and 2026 within the State Government of Santa Catarina, totaling BRL 138.76 million in approved and settled expenditures (following BRL 60.05 million retained through technical audit deductions), cross-referenced with 34,280 vehicles and 32.84 million liters of fuel. Each hypothesis was assessed against both the research-proposal wording and the pre-registered numerical threshold. In Scenario S0 (2,000 vehicles, 36 months, semi-synthetic instance), the MILP policy solved by HiGHS cost 13.66\% (BRL 6.99 million) less than the reconstructed practice while meeting all mandatory maintenance (H3 supported), although a greedy heuristic reached almost the same cost. Adding asset data to work-order features (H1) improved ROC-AUC by 0.0095 (95\% CI 0.0045--0.0146) but reduced mean absolute error by only 0.62\%, below the 15\% target (partially supported). Tree ensembles outperformed SVC-RBF in classification but had 8.65\% higher error than SVR-RBF in cost regression (H2 not supported). SHAP was exact and stable, whereas LIME showed low local fidelity (R\textsuperscript{2} = 0.24) and SHAP--LIME agreement reached a Jaccard index of 45.5\% (H4 not supported). Finally, 242 vehicles with cumulative maintenance above 40\% of book value were flagged for disposal screening. Prescriptive analytics can lower expected maintenance cost under the same budget constraints, and explanations offered to auditors should rely on methods faithful to the model.

\vspace{\onelineskip}
\noindent
\textbf{Keywords:} prescriptive analytics; public fleet management; graph theory; explainable artificial intelligence; mixed-integer linear programming.
\end{otherlanguage*}
\end{resumo}
